\scanpage{321}

\unsure{An} element of $\operatorname{Discr\acute{e}t}^{\unsure{\prime}}(\pi)$ is called a
class of discretifications of $\pi$\inserted{\unsure{or} \unsure{prediscretification} of $\pi$}.
If $\pi_0$ is a discretification, we denote its class by
$\pi_0^{\unsure{\natural}}$. The set of classes of discretifications on $\pi$
is a homogeneous \unsure{space} under $\widehat{\widehat{\unsure{\mathfrak{S}|\mathcal{G}}}}(\pi)$, the stabilizer
of $\pi_0^{\unsure{\natural}}$ \unsure{is identified with|being identified with} : the closure of
$\unsure{\mathfrak{S}|\mathcal{G}}(\pi_0)=\operatorname{Aut}_{\mathrm{lac}}(\pi_0)$ in
$\widehat{\widehat{\unsure{\mathfrak{S}|\mathcal{G}}}}(\pi)$, or again :
$\widehat{\unsure{\mathfrak{S}|\mathcal{G}}}(\pi_0)$.
The latter always contains $\pi$.
\[
\operatorname{Aut}_{\mathrm{lac}}(\pi,\pi_0^{\unsure{\natural}})
\simeq\widehat{\unsure{\mathfrak{S}|\mathcal{G}}}(\pi_0)
\]

\deleted{We have} for every discretification $\pi_0^{\unsure{\natural}}$, the subgroup of
$\widehat{\widehat{\unsure{\mathfrak{S}|\mathcal{G}}}}(\pi)$ of loop \unsure{automorphisms} which fix
$\pi_0^{\unsure{\natural}}$ is denoted $\widehat{\unsure{\mathfrak{S}|\mathcal{G}}}(\pi_0^{\unsure{\natural}})$,
it is therefore the inverse image of a subgroup of
$\widehat{\widehat{\mathcal{T}}}(\pi)$, denoted
$\widehat{\mathcal{T}}(\pi_0^{\unsure{\natural}})$. We thus have an inclusion of
extension structures

(10)
\diagram{\begin{tikzcd}[ampersand replacement=\&]
1 \arrow[r] \& \pi \arrow[r] \&
\widehat{\unsure{\mathfrak{S}|\mathcal{G}}}(\pi_0^{\unsure{\natural}}) \arrow[r] \arrow[d,hook] \&
\widehat{\mathcal{T}}(\pi_0^{\unsure{\natural}}) \arrow[r] \arrow[d,hook] \& 1 \\
1 \arrow[r] \& \pi \arrow[r] \&
\widehat{\widehat{\unsure{\mathfrak{S}|\mathcal{G}}}}(\pi) \arrow[r] \&
\widehat{\widehat{\mathcal{T}}}(\pi) \arrow[r] \& 1
\end{tikzcd}}

\margin{\uline{N.B.} The set of classes of discretifications of $\pi$ is described
\illegible{2} of the \uline{outer} group $\pi$, as
$\operatorname{Isom\,ext}(\widehat{\Pi}_{g,\nu},\pi)$\deleted{\illegible{1}}
divided by $\widehat{\mathcal{T}}_{g,\nu}$.
We moreover have \unsure{an arrow} of degree $2$
\diagram{\begin{tikzcd}[ampersand replacement=\&]
\operatorname{Isom\,ext}(\widehat{\Pi}_{g,\nu},\pi)/\widehat{\mathcal{T}}_{g,\nu}^{+}
\arrow[d] \\
\operatorname{Isom\,ext}(\widehat{\Pi}_{g,\nu},\pi)/\widehat{\mathcal{T}}_{g,\nu}
\end{tikzcd}}
which makes it possible for \unsure{every|the} \unsure{class} of discretifications, to define
its two \uline{\unsure{orientations}}, and speak of
\uline{oriented} discretifications.}

\mspage{320}

\scanpage{322}
\mspage{163}

(One shows by arithmetico-geometric means that the inclusion
$\widehat{\mathcal{T}}\longrightarrow\widehat{\widehat{\mathcal{T}}}$
is never an isomorphism.) We thus find a map
\[
\tag{11}
\operatorname{Discr\acute{e}t}^{\unsure{\prime}}(\pi)\longrightarrow
\text{Closed subgroups of }\widehat{\widehat{\mathcal{T}}}(\pi)
\]
obviously compatible with the action of
\deleted{$\unsure{\widehat{\widehat{\mathcal{T}}}}$}
$\widehat{\widehat{\unsure{\mathfrak{S}|\mathcal{G}}}}(\pi)$ (transport of structure) -- operating
on the right via $\widehat{\widehat{\mathcal{T}}}(\pi)$ and its inner automorphisms on itself.

We thus find a conjugacy class, well determined \unsure{of|of the}
closed subgroups of $\widehat{\widehat{\mathcal{T}}}(\pi)$,
which we call its ``\uline{geometric}'' Teichmüller subgroups.
\deleted{It is by no means improbable moreover that
$\widehat{\mathcal{T}}(\pi_0^{\unsure{\natural}})$ is \illegible{1}}
It could moreover be that $\widehat{\mathcal{T}}_{g,\nu}$ is
its own normalizer in $\widehat{\widehat{\mathcal{T}}}_{g,\nu}$
-- equivalent to the assertion that (11) is \unsure{injective|bijective|injection} :
The datum of a class of discretifications\inserted{of $\pi$}
would be equivalent to that of a geometric Teichmüller subgroup in
\deleted{$\unsure{\widehat{\widehat{\mathcal{T}}}}$}
its \unsure{extended} Teichmüller group $\widehat{\widehat{\mathcal{T}}}$.

\margin{this is completely screwy and ultra-false
-- a moment of aberration ! This becomes clear in what follows\ldots
The \unsure{sensible} question (with which \unsure{I} must have
\unsure{at the time} confused) \unsure{is} whether $\widehat{\mathcal{T}}$ is
\uline{invariant} in $\widehat{\widehat{\mathcal{T}}}$, i.e. whether
the normalizer is $\widehat{\widehat{\mathcal{T}}}$ \unsure{all|all} \unsure{of it} !}

\mspage{321}

\scanpage{323}

Note that we have the canonical homomorphisms
\[
\tag{12}
\widehat{\widehat{\mathcal{T}}}\deleted{\unsure{(\pi_0^{\natural})}}
\xrightarrow{\chi}\widehat{\unsure{\mathbb{Z}}}^{*}
\]
\[
\tag{13}
\widehat{\widehat{\mathcal{T}}}\longrightarrow
\unsure{\mathfrak{S}_{I}|\mathcal{G}_{I}}
\]
\margin{where $I=I(\pi)$ is \unsure{the set|the set} of conjugacy classes of loop subgroups of $\pi$}
inducing on $\widehat{\mathcal{T}}(\pi_0^{\unsure{\natural}})$
\deleted{\unsure{continuous}} corresponding homomorphisms -- whence the definition of the
subgroups $\widehat{\widehat{\mathcal{T}}}^{!}$,
$\widehat{\widehat{\mathcal{T}}}^{+}$,
$\widehat{\widehat{\mathcal{T}}}^{!+}$ and likewise for
$\widehat{\mathcal{T}}$. Note that $\chi$ takes on $\widehat{\mathcal{T}}$
only the values $\pm1$.

The subgroup $\widehat{\mathcal{T}}^{+}$ of outer automorphisms of the
profinite \unsure{groups|groups} with loops $\pi$,
\deleted{\illegible{2}}\inserted{fixing} a class of discretifications
$\pi_0^{\unsure{\natural}}$ and with multiplier $+1$, plays a very special role.
In contrast to that \unsure{which one} can conjecture about $\widehat{\mathcal{T}}$
(of which $\widehat{\mathcal{T}}^{+}$ is an \unsure{open|invariant} subgroup of index $2$),
one would rather suppose that $\widehat{\mathcal{T}}^{+}$ is invariant in
$\widehat{\widehat{\mathcal{T}}}$. In any case,
the subgroups thus defined in $\widehat{\widehat{\mathcal{T}}}$
\deleted{which}\inserted{via} \unsure{classes} of discretifications of $\pi$
\deleted{\illegible{1}} (perhaps

\mspage{322}

\scanpage{324}
\mspage{164}

there is only a single and unique one!) will be called the \uline{strict} geometric Teichmüller subgroups. The choice of such a subgroup \inserted{of $\widehat{\widehat{\mathcal{T}}}$} is, in any case, an element of structure \unsure{distinctly} weaker than that of a class of \unsure{discretifications}, and even than that of a geometric Teichmüller subgroup (not strict).

To specify the relations between these two notions, let us first recall that $\widehat{\mathcal{T}}^{+}=\Sigma$ is deduced from $\widehat{\mathcal{T}}$ as the kernel of $\chi\mid\widehat{\mathcal{T}}:\widehat{\mathcal{T}}\longrightarrow\{\pm1\}$. On the other hand, \deleted{let} (for a chosen strict geometric Teichmüller subgroup $\widehat{\mathcal{T}}^{+}$) let us consider
\[
(14)\qquad
\unsure{\mathcal{N}|N}_{\Sigma}=\mathcal{N}
=\operatorname{Norm}_{\widehat{\widehat{\mathcal{T}}}}
\bigl(\deleted{\widehat{\mathcal{T}}^{+}}\inserted{\Sigma}\bigr)
\qquad (\Sigma=\widehat{\mathcal{T}}^{+})
\]
(perhaps always equal to all of $\widehat{\widehat{\mathcal{T}}}$)

obviously (if $\deleted{\widehat{\mathcal{T}}^{+}}\inserted{\Sigma}$ comes from a $\widehat{\mathcal{T}}$)
\[
(15)\qquad
\underset{\inserted{\unsure{=}\widehat{\mathcal{T}}^{+}}}
{\deleted{\widehat{\mathcal{T}}^{+}}\inserted{\Sigma}}
\subset\widehat{\mathcal{T}}\subset\mathcal{N}_{\Sigma}.
\]
\deleted{and $\mathcal{N}/\widehat{\mathcal{T}}^{+}$ is}
The profinite group $\mathcal{N}/\deleted{\widehat{\mathcal{T}}^{+}}\inserted{\Sigma}$
\deleted{is} associated with \deleted{the \unsure{structure}}
$\deleted{\widehat{\mathcal{T}}^{+}}\inserted{\Sigma}$ is denoted
$\boldsymbol{\Gamma}_{\deleted{\unsure{\widehat{\mathcal{T}}^{+}}}\inserted{\Sigma}}$
\deleted{\illegible{1}} \deleted{\illegible{2}}
(if $\deleted{\widehat{\mathcal{T}}^{+}}\inserted{\Sigma}$ is unique, we write $\boldsymbol{\Gamma}_{\pi}$);

\mspage{323}

\scanpage{325}

and the $\widehat{\mathcal{T}}\deleted{\unsure{\mathcal G}}$ \unsure{giving} rise to the same $\deleted{\widehat{\mathcal{T}}^{+}}\inserted{\Sigma}$ correspond \deleted{\unsure{to the classes}} to a conjugacy class of elements of order $2$ of $\boldsymbol{\Gamma}_{\deleted{\unsure{\widehat{\mathcal{T}}^{+}}}\inserted{\Sigma}}$. It is clear \deleted{\unsure{that we|that}} in any case that we have an injective map
\[
(16)\qquad
\begin{gathered}
\text{set of geometric}\\
\text{Teichmüller subgroups}\\
\widehat{\mathcal{T}}\text{ of }\widehat{\widehat{\mathcal{T}}}\text{ such that}\\
\widehat{\mathcal{T}}^{+}=\Sigma
\end{gathered}
\longrightarrow
\begin{gathered}
\text{set of }\unsure{\text{elements}|\text{elements}}\\
\text{of order }2\text{ of }\boldsymbol{\Gamma}_{\Sigma}
\end{gathered}
\quad :
\]
and that its image is stable under conjugation. Let us show that \unsure{two|the} elements of the image are conjugate under $\boldsymbol{\Gamma}_{\Sigma}$. Indeed, let
\[
\widehat{\mathcal{T}},\ \widehat{\mathcal{T}}^{\prime}\supset\Sigma
\quad\text{such that}\quad
\Sigma=\widehat{\mathcal{T}}^{+}=\widehat{\mathcal{T}}^{\prime+}.
\]
\deleted{We \illegible{1} \unsure{obviously}}
\inserted{$\exists\,g\in\widehat{\widehat{\mathcal{T}}}$}
such that
$\widehat{\mathcal{T}}^{\prime}=\operatorname{int}(g)\unsure{\mathcal{T}}$,
and we \unsure{will have|have} therefore
$\widehat{\mathcal{T}}^{\prime+}=\operatorname{int}_{g}\widehat{\mathcal{T}}^{+}$
i.e. $g\in\mathcal{N}$, OK.

The elements of order $2$ thus obtained in $\boldsymbol{\Gamma}_{\Sigma}$ are called the canonical involutions. We will give a remarkable conjectural interpretation of them below. The set $\widehat{\widehat{\mathcal{T}}}/\mathcal{N}$ is identified with the set of \unsure{subgroups} of

\mspage{324}

\scanpage{326}
\mspage{165}

strict geometric Teichmüller (once the ``origin'' element $\Sigma$ has been chosen). More \unsure{intrinsically}, we \unsure{will have}:
\[
(17)\qquad
\begin{gathered}
\text{Set of strict geometric}\\
\text{Teichmüller subgroups}\\
\text{of }\widehat{\widehat{\mathcal{T}}}(\pi)
\end{gathered}
\unsure{\xleftarrow{\sim}|\simeq}
\operatorname{Isomext}_{\unsure{\mathrm{lac}}}
(\widehat{\Pi}_{g,\nu},\pi)/\unsure{\mathcal{N}_{g,\nu}}
\]
where we set, as is only right
\[
(18)\qquad
\mathcal{N}_{g,\nu}
=\operatorname{Norm}_{\widehat{\widehat{\mathcal{T}}}_{g,\nu}}
\bigl(\widehat{\mathcal{T}}_{g,\nu}^{+}\bigr)
\]
(perhaps equal to all of $\widehat{\widehat{\mathcal{T}}}_{g,\nu}$!)

Let us now consider the \uline{modular topos} \uline{over} $\mathbb{Q}$ of algebraic curves of type $(g,\nu)$, denoted $M_{g,\nu\,\mathbb{Q}}$, or simply $M_{g,\nu}$. If we choose a universal covering $\widetilde{M}_{g,\nu}$ of $M_{g,\nu\,\mathbb{Q}}$ (hence a universal covering of $\unsure{\operatorname{Spec}}\mathbb{Q}$ i.e. an algebraic closure $\overline{\mathbb{Q}}$ of $\mathbb{Q}$) we can \unsure{specify} the $\pi_1(M_{g,\nu\,\mathbb{Q}})$ as the group of $M_{g,\nu\,\mathbb{Q}}$-\unsure{automorphisms} of this \unsure{covering|covering}; \deleted{\unsure{and of it} \unsure{make} \illegible{1}} \inserted{call it the} \uline{arithmetic Teichmüller group} of type $(g,\nu)$ (relative to the choice of $\widetilde{M}_{g,\nu\,\mathbb{Q}}$). We will thus have an exact sequence

\mspage{325}

\scanpage{327}

\[
(19)\qquad 1\longrightarrow \pi_1(M_{g,\nu\,\overline{Q}})
\longrightarrow \pi_1(M_{g,\nu})
\longrightarrow \operatorname{Gal}(\overline{Q}/Q)\longrightarrow 1
\]
\deleted{\(1\longrightarrow \unsure{\widehat{\mathcal{T}}_{g,\nu}}\longrightarrow \unsure{\Gamma}\)}

(\unsure{where} $\overline{Q}$ and the base points are \unsure{made explicit} as has been said) Let us note that by the \unsure{transcendental} theory of Teichmüller, we have an isomorphism (defined modulo \unsure{operation} \unsure{induced} by $\pi_1(M_{g,\nu})$ on $\pi_1(M_{g,\nu,\overline{Q}})$)
\[
(20)\qquad \pi_1(\unsure{M_{g,\nu,\overline{Q}}})\simeq\widehat{\mathcal{T}}_{g,\nu}^{+}
\qquad (\text{profinite compactification of }\mathcal{T}_{g,\nu}^{+})
\]

Let us on the other hand consider the \unsure{scheme.} $\unsure{U|\mathcal{U}}_{g,\nu}$\inserted{ over $M_{g,\nu}$}, ``universal'' curve of type $(g,\nu)$

We thus have, on choosing a universal covering $\widetilde{\unsure{U|\mathcal{U}}}_{g,\nu}$ of it \uline{\unsure{over}} the chosen universal covering $\unsure{\widetilde{M}_{g,\nu}^{\natural}}$ of $M_{g,\nu}$, an exact sequence
\diagram{
\begin{tikzcd}[column sep=large,row sep=large,ampersand replacement=\&]
(21) \& 1 \arrow[r] \& \pi \arrow[r] \arrow[dd,"\Vert" description] \& \pi_1(\unsure{\overset{\deleted{\sim}}{\unsure{U|\mathcal{U}}}_{g,\nu}}) \arrow[r] \arrow[dd] \& \pi_1(M_{g,\nu}) \arrow[r] \arrow[dd] \& 1 \\
\text{which maps to:} \& \& \& \& \& \\
(22) \& 1 \arrow[r] \& \pi \arrow[r] \& \widehat{\widehat{\unsure{\mathfrak{S}|\mathcal{G}}}}(\pi) \arrow[r] \& \widehat{\widehat{\mathcal{T}}}(\pi) \arrow[r] \& 1
\end{tikzcd}
}
and the same for $\unsure{U|\mathcal{U}}_{g,\nu\,\overline{Q}}\longrightarrow M_{g,\nu\,\unsure{\overline{Q}|Q}}$

\mspage{326}

\scanpage{328}
\mspage{\unsure{166|168}}

\[
\unsure{(\overline{21})|(\overline{21}')}\qquad
1\longrightarrow\pi\longrightarrow\pi_1(\unsure{U|\mathcal{U}}_{g,\nu\,\overline{Q}})
\longrightarrow\pi_1(M_{g,\nu\,\overline{Q}})\longrightarrow1
\]
\deleted{\(\unsure{(\overline{22})}\qquad 1\longrightarrow\pi\longrightarrow\)}

We of course have an \unsure{isomorphism} of groups with \unsure{loops|ends}
\[
(23)\qquad \pi\simeq\widehat{\Pi}_{g,\nu}
\]
whose indeterminacy we would like to determine precisely

Let $M_{g,\nu\,C}$ denote the complex modular Teichmüller \deleted{\illegible{1}}\inserted{topos}, defined using the \deleted{differentiable}\inserted{$C^\infty$} surface $\unsure{U|\mathcal{U}}_{g,\nu}$, it is equipped with the canonical universal covering $\widetilde{M}_{g,\nu\,C}$ (\S~21)
\margin{$\widetilde{M}_{g,\nu\,C}=E_g/\unsure{A_{g,\nu}^{0}|A_{g,\nu}^{\circ}}$}
with group \unsure{of automorphisms} $\mathcal{T}_{g,\nu}^{+}$ precisely \unsure{--|...}
$\unsure{U|\mathcal{U}}_{g,\nu\,C}$ being itself also equipped with a universal covering \unsure{over} the preceding one,
$\widetilde{\unsure{U|\mathcal{U}}}_{g,\nu\,C}=E_g/\unsure{A_{g,\nu+1}^{0}|A_{g,\nu+1}^{\circ}}$.
This gives rise to the exact sequence of discrete groups (in the general topology context)
\diagram{
\begin{tikzcd}[column sep=large,row sep=large,ampersand replacement=\&]
(24) \& 1 \arrow[r] \& \Pi_{g,\nu} \arrow[r] \arrow[dd] \& \unsure{\mathfrak{S}|\mathcal{G}}_{g,\nu}^{+} \arrow[r] \arrow[d,equal] \& \mathcal{T}_{g,\nu}^{+} \arrow[r] \arrow[d,equal] \& 1 \\
\& \& \& \pi_1(\unsure{U|\mathcal{U}}_{g,\nu\,C}) \arrow[d] \& \pi_1(M_{g,\nu\,C}) \arrow[d] \& \\
\text{(mapping to} \& 1 \arrow[r] \& \Pi_{g,\nu} \arrow[r] \& \operatorname{Aut}_{\mathrm{lac}}\deleted{\unsure{\mathrm{ext}}}(\Pi_{g,\nu}) \arrow[r] \arrow[d,equal] \& \operatorname{Aut\,ext}_{\mathrm{lac}}(\Pi_{g,\nu}) \arrow[r] \arrow[d,equal] \& 1 \\
\& \& \& \unsure{\mathfrak{S}|\mathcal{G}}_{g,\nu} \& \mathcal{T}_{g,\nu} \&
\end{tikzcd}
}
by an \unsure{isomorphism|homomorphism} of extension \unsure{structures}), which by passage to the profinite completions gives the analogous exact sequence for the

\mspage{327}

\scanpage{329}

algebraic multiplicities over $\mathbb C$

\diagram{\begin{tikzcd}[ampersand replacement=\&,column sep=1.3em,row sep=1.2em]
(25) \& 1 \arrow[r] \&
\widehat{\Pi}_{g,\nu} \arrow[r] \arrow[d,equal] \&
\deleted{\unsure{\widehat{\Pi}}}\,\widehat{\mathcal G}_{g,\nu}
  \arrow[r] \arrow[d,very thick] \&
\widehat{\mathcal T}_{g,\nu}^{+} \arrow[r]
  \arrow[d,phantom,"{\rotatebox{-90}{$\unsure{\hookrightarrow}$}}" description] \& 1 \\
\& \& \inserted{\widehat{\Pi}_{g,\nu}} \arrow[r] \&
\inserted{\operatorname{Aut}\deleted{\unsure{\mathrm{ext}}}_{\mathrm{lac}}(\widehat{\Pi}_{g,\nu})}
  \arrow[r] \&
\inserted{\operatorname{Autext}_{\mathrm{lac}}(\widehat{\Pi}_{g,\nu})}
  \& \\[-0.7em]
\& \& \& |[xshift=-2.2em]| {\inserted{\raisebox{1ex}{\rotatebox[origin=c]{-60}{$=$}}\,\unsure{\widehat{\widehat{\mathcal G}}}_{g,\nu}}} \&
|[xshift=-2.2em]| {\inserted{\raisebox{1ex}{\rotatebox[origin=c]{-60}{$=$}}\,\widehat{\widehat{\mathcal T}}_{g,\nu}}} \&
\end{tikzcd}}

This situation moreover \unsure{arises} \unsure{canonically} from an analogous situation over $\overline{\mathbb Q}_0=\text{algebraic closure of }\mathbb Q\text{ in }\mathbb C$, the choice of $\widetilde{\unsure{U|\mathcal{U}}}_{g,\nu\mathbb C}$ defining a $\widetilde{\unsure{U|\mathcal{U}}}_{g,\nu\overline{\mathbb Q}_0}$, hence a $\widetilde{M}_{g,\nu\overline{\mathbb Q}_0}$.

To define an isomorphism between this sequence (25) and the exact sequence $(\overline{21})$ \deleted{one must} (\unsure{equipped} with \unsure{$2$ homomorphisms|its homomorphism} to (22)) one must

$1^{\circ})$ choose an isomorphism $\overline{\mathbb Q}\simeq\overline{\mathbb Q}_0$ --- the indeterminacy is in $\mathrm{Gal}(\overline{\mathbb Q}/\mathbb Q)$ ; which makes it possible to identify $\overline{\mathbb Q}$ and $\overline{\mathbb Q}_0$ ;

$2^{\circ})$ Choose an isomorphism (\deleted{\unsure{between}} $\overline{\mathbb Q}=\overline{\mathbb Q}_0$) of $(\widetilde{\unsure{U|\mathcal{U}}}_{g,\nu})_0$ (constructed by \unsure{transcendental|transcendental} means over $\mathbb C$ and descended to $\overline{\mathbb Q}_0$) with $\widetilde{\unsure{U|\mathcal{U}}}_{g,\nu\overline{\mathbb Q}}$ --- the indeterminacy is in $\pi_1(M_{g,\nu\overline{\mathbb Q}})$.

In summary, we have a class of isomorphisms of \unsure{(24)|(21)}\inserted{\unsure{2|?}} and (25), defined modulo automorphisms \inserted{\deleted{\illegible{1}} \illegible{1} \unsure{once we have chosen} $\overline{\mathbb Q}\simeq\overline{\mathbb Q}_0$} inner in \inserted{the \unsure{\uline{profinite}} group} $\pi_1(\unsure{U|\mathcal{U}}_{g,\nu})$. (\unsure{These|The} \inserted{\unsure{\uline{corresponding}}} isomorphisms transform the \deleted{structure} \unsure{Standard|standard} \uline{oriented} \uline{class of \unsure{discretification}} of $\widehat{\Pi}_{g,\nu}$ into \deleted{\unsure{a class}} \inserted{exactly} \uline{one} class of oriented discretifications of $\pi$.

\mspage{328}

\scanpage{330}
\mspage{167}

\deleted{Thus $\pi$ of (21) is \unsure{canonically} equipped with an oriented \deleted{\unsure{class of}} \inserted{\unsure{quasi}} discretification \inserted{\unsure{$\pi_0^{\natural}$}}, and in the homomorphism from $(\overline{21})$ to (22), we find that $\pi_1(\unsure{U|\mathcal{U}}_{g,\nu\overline{\mathbb Q}})$ is identified with $\widehat{\unsure{\mathfrak{S}|\mathcal{G}}}(\unsure{\pi_0^{\natural}})$, and $\pi_1(M_{g,\nu})$ with $\widehat{\mathcal T}(\unsure{\pi_0^{\natural}})$ \unsure{$\longrightarrow$}}

Changing the choice of the isomorphism to its complex conjugate, we find the opposite oriented \inserted{\unsure{quasi-}} \unsure{discretification}.

We see on the other hand that the image of $\pi_1(M_{g,\nu\overline{\mathbb Q}})\unsure{\subset}\pi_1(M_{g,\nu})$ in $\widehat{\widehat{\mathcal T}}(\pi)$ is none other than the $\widehat{\mathcal T}^{+}$ \unsure{corresponding|corresponding} to an \inserted{arbitrary} one of the chosen classes of discrete quasi-rigidifications --- the subgroup does not depend on the choice of this class --- which merely expresses that the image of $\pi_1(M_{g,\nu\overline{\mathbb Q}})$ in $\widehat{\widehat{\mathcal T}}(\pi)$, \unsure{considered|considered} as a subgroup of the image of the larger group $\pi_1(M_{g,\nu})$, is \unsure{normal|normal} in it --- which follows from the fact that $\pi_1(M_{g,\nu\overline{\mathbb Q}})$ is \unsure{normal} in $\pi_1(M_{g,\nu\unsure{\mathbb Q|\overline{\mathbb Q}}})$. Thus $\pi$

\mspage{329}
